
\subsubsection{5.1 Corpus Coverage (RQ1)}

\textbf{Table 2: Collection and Coverage Results by Repository}

\begin{table*}[htbp]
\centering
\resizebox{\dimexpr 2\columnwidth+\columnsep\relax}{!}{%
\begin{tabular}{lllll}
\toprule
Repository & Collected & Scoreable & Coverage Rate & Gate (>= 0.70) \\
\midrule
HuggingFace Hub & 496 & 496 & 1.000 & Pass \\
OpenML & 200 & 200 & 1.000 & Pass \\
UCI ML Repository & 62 & 0 & 0.000 & Fail \\
**Total** & **758** & **696** & **0.918** & **Pass** \\
\bottomrule
\end{tabular}
}
\end{table*}

The overall coverage rate of 0.918 exceeds the 0.70 threshold. HuggingFace Hub and OpenML each achieve 100\% coverage, indicating that their structured APIs contain sufficient fields for DTS scoring.

The UCI result requires explanation: all 62 collected datasets returned null values for all DTS-mapped fields. This is attributable to a field naming mismatch between the DTS scorer's configured UCI field keys and the actual key names returned by the \texttt{ucimlrepo} library. The UCI datasets likely contain some relevant metadata under different field names, but the current field mapping does not capture it. This is an engineering issue in the pipeline configuration, not evidence that UCI datasets lack all documentation. Excluding UCI from the denominator yields 100\% coverage for HuggingFace and OpenML combined.

\begin{figure}[htbp]
\centering
\includegraphics[width=\columnwidth]{figures/gate_metrics_comparison.png}
\caption{Gate metrics comparison}
\label{fig:gate_metrics_comparison}
\end{figure}

\textit{Figure 1: Gate criteria achievement. Both pre-specified thresholds (coverage >= 0.70 and proxy r >= 0.70) are met.}

\subsubsection{5.2 Internal Consistency of DTS Weighting (RQ2)}

\textbf{Table 3: Gate Criteria Results}

\begin{table*}[htbp]
\centering
\resizebox{\dimexpr 2\columnwidth+\columnsep\relax}{!}{%
\begin{tabular}{lllll}
\toprule
Criterion & Threshold & Achieved & Margin & Status \\
\midrule
Coverage rate & >= 0.70 & 0.918 & +21.8 pp & Pass \\
Proxy Pearson r & >= 0.70 & 0.989 & +28.9 pp & Pass \\
\bottomrule
\end{tabular}
}
\end{table*}

The mean weighted DTS score is 0.169 (SD = 0.124), compared to a mean unweighted score of 0.229 (SD = 0.150). The weighted score is approximately 26\% lower than the unweighted score ((0.229 - 0.169) / 0.229 = 0.262). This difference arises because the highest-weighted DTS sections (Preprocessing, Uses) are rarely populated in API responses, pulling the weighted average down relative to naive field counting.

The proxy Pearson correlation between weighted and unweighted scores is r = 0.989 (p = 5.77 x $10^-101$, 95\% bootstrap CI [0.985, 0.994], n = 120). This near-perfect correlation is partially expected by construction: both measures are computed from the same binary field presence values and differ only in the weight vector. The correlation confirms that the scoring algorithm behaves consistently, but it does not constitute external validation against independent human expert judgment. The primary distinction between the two scoring approaches is in score magnitude (0.169 vs. 0.229), not in dataset rank ordering.

\begin{figure}[htbp]
\centering
\includegraphics[width=\columnwidth]{figures/dts_score_distribution.png}
\caption{DTS score distribution}
\label{fig:dts_score_distribution}
\end{figure}

\textit{Figure 2: Distribution of DTS-weighted and unweighted scores across the 758-dataset corpus.}

\begin{figure}[htbp]
\centering
\includegraphics[width=\columnwidth]{figures/human_automated_scatter.png}
\caption{Proxy validation scatter}
\label{fig:human_automated_scatter}
\end{figure}

\textit{Figure 3: Scatter plot of weighted vs. unweighted DTS scores on the 120-dataset validation subsample, with bootstrap confidence interval annotation.}

\subsubsection{5.3 Per-Section Coverage Profile (RQ3)}

\textbf{Table 4: Per-Section DTS Coverage Rates by Repository}

\begin{table*}[htbp]
\centering
\resizebox{\dimexpr 2\columnwidth+\columnsep\relax}{!}{%
\begin{tabular}{llllll}
\toprule
DTS Section & DTS Weight & Overall & HuggingFace & OpenML & UCI \\
\midrule
Motivation & 1.0 & 0.547 & 0.647 & 0.470 & 0.000 \\
Composition & 0.9 & 0.267 & 0.105 & 0.750 & 0.000 \\
Collection & 2.1 & 0.184 & 0.147 & 0.333 & 0.000 \\
Preprocessing & 1.8 & 0.002 & 0.000 & 0.008 & 0.000 \\
Uses & 1.5 & 0.000 & 0.000 & 0.000 & 0.000 \\
Distribution & 0.7 & 0.247 & 0.190 & 0.465 & 0.000 \\
\bottomrule
\end{tabular}
}
\end{table*}

\begin{figure}[htbp]
\centering
\includegraphics[width=\columnwidth]{figures/per_section_coverage_heatmap.png}
\caption{Per-section coverage heatmap}
\label{fig:per_section_coverage_heatmap}
\end{figure}

\textit{Figure 4: Per-section DTS coverage rates across repositories. Preprocessing (weight 1.8) and Uses (weight 1.5) score near-zero across all repositories. These sections exist as free-text prose in dataset cards rather than structured API fields, which is why automated API-based scoring cannot detect them.}

The two highest-weighted DTS sections have near-zero coverage universally:

\begin{itemize}
\item Preprocessing (weight 1.8): overall coverage 0.002. Only OpenML shows any non-zero value (0.008), corresponding to a small number of datasets with a structured preprocessing field.
\item Uses (weight 1.5): overall coverage 0.000. No dataset in the 758-dataset corpus documented intended uses or known limitations through structured API fields.
\end{itemize}

The sections with the highest coverage -- Motivation (0.547), Composition (0.267), Distribution (0.247) -- have the lowest DTS weights (1.0, 0.9, 0.7 respectively). Collection (weight 2.1) achieves 0.184 overall coverage, indicating partial availability of structured collection-related fields.

This pattern arises because preprocessing steps, intended uses, and known limitations are typically described in free-text prose within dataset cards and README files. HuggingFace's \texttt{card\textbackslash \_data} API returns only structured YAML key-value pairs; the prose content of dataset cards is not indexed as structured fields. This structural mismatch between where documentation exists (prose) and what the API exposes (structured fields) is the documentation API gap.

\begin{figure}[htbp]
\centering
\includegraphics[width=\columnwidth]{figures/missing_field_analysis.png}
\caption{Missing field analysis}
\label{fig:missing_field_analysis}
\end{figure}

\textit{Figure 5: Field-level missing rates by DTS section and repository, providing granular detail complementing the section-level heatmap.}

\subsubsection{5.4 Directional Comparison with Prior Manual Scoring}

The per-section pattern observed here -- higher coverage for lower-weighted sections and near-zero coverage for Preprocessing and Uses -- is directionally consistent with the patterns described in Rondina et al. (2025), who reported near-zero Preprocessing and Collection coverage and higher Motivation coverage in their manual 100-dataset study. However, exact numerical comparison is not possible because (a) Rondina et al. scored manually while we score from API fields, measuring partially different constructs, and (b) the Rondina et al. citation has not been independently verified. The directional consistency suggests that the section asymmetry is structural rather than an artifact of either scoring method.

\subsubsection{5.5 Mechanism Activation Summary}

\textbf{Table 5: Mechanism Activation Indicators}

\begin{table}[htbp]
\centering
\resizebox{\columnwidth}{!}{%
\begin{tabular}{lll}
\toprule
Indicator & Expected & Observed \\
\midrule
Scoring pipeline executes without errors & True & True \\
Coverage achievable for at least 2 repositories & True & True (HF: 1.000, OpenML: 1.000) \\
DTS weighting produces distinct score level & weighted != unweighted & True (0.169 vs. 0.229) \\
Positive proxy correlation & r > 0 & True (r = 0.989) \\
\bottomrule
\end{tabular}
}
\end{table}
